\chapter{Sucesiones y series de funciones}
\label{cap:sucesiones-funciones}

La convergencia de funciones requiere controlar simultáneamente todos los puntos del dominio. Se comparan convergencia puntual y uniforme, y se estudian condiciones que permiten conservar continuidad e intercambiar límites con integrales o derivadas \cite{abbott2015,rudin1976}.

\section{Convergencia puntual}
Definición punto a punto y ejemplos en los que cambia la regularidad del límite.

\section{Convergencia uniforme}
Definición uniforme, criterios y comparación con convergencia puntual.

\section{Criterio de Cauchy uniforme}
Criterio de Cauchy para sucesiones de funciones y completitud del espacio de funciones acotadas.

\section{Preservación de la continuidad}
Condiciones bajo las cuales el límite uniforme de funciones continuas es continuo.

\section{Intercambio de límite e integral}
Hipótesis que justifican pasar el límite bajo el signo integral.

\section{Intercambio de límite y derivada}
Convergencia de derivadas y condiciones para derivar el límite.

\section{Series de funciones}
Convergencia puntual y uniforme de series y convergencia normal.

\section{Criterio \texorpdfstring{$M$}{M} de Weierstrass}
Prueba de convergencia uniforme a partir de cotas sumables independientes del punto.
